\subsection{Universal tensor inversion of the actual coefficient}

The finite-dimensional inverse obstruction in
\fid{F38}/\fid{F62} has a stronger categorical form.
Throughout, an additive monoidal category has finite
biproducts and tensor product additive in both variables.
It is nonzero when the identity of its unit is nonzero.
If \(1_1=0\), the unit constraint identifies
\(1_Y\) with \(1_1\otimes1_Y=0\), so every object is zero.
A tensor inverse is two-sided. A strong monoidal
functor has invertible tensor and unit constraints;
we also require it to be additive.

\begin{finding}[\fid{F68}, universal coefficient-inversion obstruction]
\label{f:universalCoefficientInversion}
Let \(F:\mathcal C\to\mathcal D\) be an additive
strong monoidal functor to a nonzero additive
monoidal category.

If an object of \(\mathcal C\) is a finite direct
sum of at least two tensor-invertible objects,
its image under \(F\) cannot be tensor-invertible.

If \(\mathcal C\) is symmetric and \(X\) is
dualizable with categorical dimension zero, then
\(F(X)\) likewise cannot be invertible.
This second assertion does not require a target
symmetry or require \(F\) to preserve one.

Consequently the actual torus compact-cohomology
object, its stated positive powers and Tate/shift
normalizations, and the actual projective-line
cohomology object cannot be inverted by such a functor.
The relative coefficient \(\mathcal K\) of
\fid{F62} is also excluded when its actual dual
pair and source symmetry maps are retained.
No boundedness or faithful-realization hypothesis
on \(\mathcal D\) is needed.
\end{finding}

\begin{proof}
\par\smallskip\noindent\textbf{An invertible summand of the unit.}\quad
Suppose \(U\) is invertible and a summand of the unit.
Choose \(i:U\to1\), \(p:1\to U\) with \(pi=1_U\),
and put \(e=ip\).
For every object \(Y\), the endomorphism
\[
 e_Y=\lambda_Y(e\otimes1_Y)\lambda_Y^{-1}
\]
is natural in \(Y\), where \(\lambda\) is the
left unitor. Coherence gives \(e_1=e\).
Naturality with \(i\) gives \(ei=i e_U\).
Since \(ei=i\), composition with \(p\) gives
\(e_U=1_U\).
Right tensoring with invertible \(U\) is an
equivalence, hence is faithful on
\(\operatorname{End}(1)\). Thus \(e=1_1\).
If \(1=U\oplus V\), its projection onto \(U\)
therefore equals the identity, and the inclusion
of \(V\) is zero. Hence \(V=0\).
No braiding was used.

Now suppose \(P=\bigoplus_{i=1}^rL_i\), \(r\ge2\),
with each \(L_i\) invertible, were itself invertible.
Right tensor its biproduct by an inverse.
The unit becomes a sum of invertible objects
\(L_i\otimes P^{-1}\).
The preceding argument kills the complement of
the first summand, hence kills another invertible
summand. Tensoring that zero object with its inverse
makes the unit zero, a contradiction.
An additive strong monoidal functor preserves
the displayed biproduct and all the invertibility
data, proving the first assertion.

\par\smallskip\noindent\textbf{The zero source-trace lemma.}\quad
Choose a dual pair in \(\mathcal C\),
\[
 \eta:1\to X\otimes X^\vee,\qquad
 \epsilon:X^\vee\otimes X\to1,
\]
with its triangle identities. The actual
source dimension endomorphism is
\begin{equation}\label{eq:sourceZeroTraceWord}
 \dim_{\mathcal C}(X)
       =\epsilon\,c_{X,X^\vee}\,\eta=0,
\end{equation}
where \(c\) is the source symmetry.
Strong monoidality transports this dual pair,
including its triangle identities, to
\(F(X),F(X^\vee)\).

If an object \(Z\) is invertible, any dual
\(Z^\vee\) is its tensor inverse and both
evaluation and coevaluation are isomorphisms.
Indeed the dual pair supplies
\[
 (-\otimes Z)\dashv(-\otimes Z^\vee).
\]
The left functor is an equivalence, so its right
adjoint is an inverse equivalence. The adjunction
unit and counit are isomorphisms; at the unit
object they are the displayed coevaluation and
evaluation, with the monoidal constraints.

Thus, if \(F(X)\) were invertible, the transported
\(F(\eta)\) and \(F(\epsilon)\) would be
isomorphisms. So is \(F(c)\), since \(c\) is one.
Their composite is the image of
\eqref{eq:sourceZeroTraceWord}, hence is zero
by additivity. A zero endomorphism of the unit
that is also an isomorphism forces \(1_1=0\).
This proves the second assertion.
We have not identified \(F(c)\) with a target
braiding. The actual source dual pair and source
symmetry morphism are the retained data.

When \(F\) is symmetric, the familiar shorter
argument is also available: categorical dimension
is preserved and the dimensions of an invertible
object and its inverse multiply to \(1_1\).
The source-trace-word proof shows why target
symmetry is unnecessary.

\par\smallskip\noindent\textbf{An actual derived Weil splitting of the torus.}\quad
Fix \(\mathbb F_q\), \(\ell\nmid q\), and a
finite extension \(E/\mathbb Q_\ell\).
Write \(E(r)\) for the Tate line with geometric
Frobenius \(q^{-r}\). The cohomology objects below
are geometric cohomology with their Weil action.
Put
\[
 T=R\Gamma_c(\mathbb G_{m,\overline{\mathbb F}_q},E).
\]
The smooth-pair Gysin sequence for
\(\{0,\infty\}\subset\mathbb P^1\) gives
\[
 H^0(\mathbb G_m,E)=E,\qquad
 H^1(\mathbb G_m,E)=E(-1),
\]
with no other cohomology. The map
\(E(-1)^2\to H^2(\mathbb P^1,E)=E(-1)\)
sends the two rational point classes to the
degree-one generator, so is the sum map
\cite[Theorem~16.1 and Corollary~16.2, p.~108;
Theorem~23.2, pp.~139--140;
Remark~24.2(b), p.~145]{Milne}.

The rational point \(1\in\mathbb G_m(\mathbb F_q)\)
gives actual Frobenius-equivariant derived maps
\[
 E\xrightarrow{\mathrm{unit}}R\Gamma(\mathbb G_m,E)
     \xrightarrow{\mathrm{ev}_1}E
\]
whose composite is \(1_E\).
Thus \(E\) splits off as an actual derived object.
The complement has only \(E(-1)\) in degree one,
so truncation identifies it with \(E(-1)[-1]\):
\[
 R\Gamma(\mathbb G_m,E)\simeq E\oplus E(-1)[-1].
\]
No Frobenius-semisimple splitting has been assumed.
The canonical cup-product/trace pairing for the
smooth curve, perfect on cohomology, gives an
actual derived Weil isomorphism
\[
 T\simeq R\Gamma(\mathbb G_m,E)^\vee(-1)[-2]
 \simeq E[-1]\oplus E(-1)[-2]
\]
\cite[Theorem~24.1(a),(b), p.~144]{Milne}.
The trace Tate line and naturality make this
isomorphism Frobenius-equivariant.

Both displayed summands are invertible.
Their categorical dimensions under the usual
Koszul symmetry are \(-1,+1\), so
\[
 \dim(T)=0,\qquad
 \operatorname{Tr}_{\mathrm{Eul}}(F^n;T)=q^n-1.
\]
The second, inserted Frobenius trace is distinct
from the first, categorical trace of the identity.
Both categorical lemmas now exclude inversion
of \(T\) by the stated functors.

K\"unneth preserves the actual decompositions,
so positive tensor powers and their Tate/shift
normalizations are excluded too
\cite[Theorem~22.1 and Remark~22.3(b)--(d),
pp.~135--136, with the following \(\ell\)-adic passage]{Milne}.
For the \fid{F62}/IV coefficient specifically,
\[
 T^{\otimes2}(1)\simeq
 E(1)[-2]\oplus E[-3]^{\oplus2}\oplus E(-1)[-4].
\]
Its trace is \(q^n+q^{-n}-2\).
The exact VIII coefficient object
\(R\Gamma_c(\mathbb G_m^2\times\mathbb P^2,E(2))[1]\)
similarly splits into twelve shifted Tate
invertibles: use K\"unneth and the projective
hyperplane maps giving the three summands of
\(R\Gamma(\mathbb P^2,E)\).

\par\smallskip\noindent\textbf{The projective-line object in \fid{F38}.}\quad
The rational-point retraction and projective-line
cohomology give the actual derived Weil splitting
\[
 H=R\Gamma(\mathbb P^1,E)
       \simeq E\oplus E(-1)[-2].
\]
Here \(\dim(H)=2\); the zero-trace lemma is not
being applied. The invertible-summand lemma
nevertheless excludes making \(H\) invertible
by any additive strong monoidal functor to a
nonzero target.
This strengthens the earlier faithful-realization
boundary. A stack coefficient with inverse
Euler traces still need not be a tensor inverse
of the actual object \(H\).

\par\smallskip\noindent\textbf{The relative coefficient without a global splitting.}\quad
Retain the actual \fid{F62} base
\(O=(\mathbb P^1\times\mathbb P^1)\setminus\Delta\),
line bundle \(\mathcal L=O(1,-1)|_O\), and
\[
 \mathcal K=R\pi_!E,\qquad
 \pi:\mathcal L^\times\to O.
\]
Its established triangle is
\[
 \mathcal K\longrightarrow E_O(-1)[-2]
   \xrightarrow{c_1(\mathcal L)}E_O
   \longrightarrow\mathcal K[1].
\]
The Chern class is nonzero. We do not assert
the global splitting into the two torus summands.

On a Zariski cover trivializing \(\mathcal L\),
compact-support base change and K\"unneth identify
\(\mathcal K\) with the constant complex of fiber \(T\).
Thus it is locally a bounded complex of finite-rank
free \(E\)-sheaves. It is tensor-dualizable, with
dual \(R\mathcal Hom(\mathcal K,E_O)\);
the evaluation identities can be checked on this cover.
Its categorical identity trace restricts to zero
on every trivializing open. It is a degree-zero
endomorphism of \(E_O\), so it is zero globally.

The zero-source-trace lemma now excludes inversion
of \(\mathcal K\) by a functor retaining its
dual pair and source symmetry maps. The target
need not be symmetric.
This uses dualizability and the actual zero
source trace, not a nonexistent global direct sum.
The nonzero Chern extension is preserved.
The relative two-torus coefficient and its stated
projective/Tate normalizations have zero categorical
trace for the same local reason and are excluded too.
\end{proof}

\paragraph{What the universal statement does and does not invert.}
An additive strong monoidal localization making
one of these actual objects invertible is the
zero category. Exact strong monoidal functors
to stable or unbounded targets are included;
the proof only uses finitely many source objects
and maps.
By contrast, a Grothendieck or trace ring may
formally invert the corresponding Laurent factor.
For example,
\(\mathbb Z[\mathbb L,\mathbb L^{-1},
(\mathbb L-1)^{-1}]\) is a nonzero ring.
That does not invert the actual biproduct
\(E[-1]\oplus E(-1)[-2]\) in the categorical sense.

One-sided inverses, nonadditive or lax monoidal
operations, a different object having the same trace,
and arbitrary filtered extraction are outside
these assertions.
The trace argument needs the specified source
dual pair and symmetry morphism.
The split-summand argument only needs its biproduct
and the invertibility data of the summands; merely
removing duality maps does not evade that argument
if its own data remain.
No assertion is made that every normalization
or canonical coefficient extraction must invert
the actual torus object by such a functor.

% Separate paragraph for Part III.
\paragraph{The \fid{F68} distinction between actual duality and virtual inversion.}
The bounded coefficient objects are perfect over \(E\)
in SGA~6, I, Definition~4.7, pp.~121--122 and
Example~4.8, p.~122 \cite{SGA6}.
Their direct sums, dual pairs and symmetry morphisms
are actual categorical data.
For the relative line bundle, EGA~I,
Theorem~1.3.7, p.~86, Corollary~1.3.8, p.~87,
and Theorem~1.4.1, pp.~90--92 supply the
affine module/sheaf dictionary \cite{EGA}.
The exact geometric Gysin, projective-bundle,
K\"unneth and duality locators were cited in the proof.

SGA~5, III, 4.7 is the existing trace-formula
framework anchor \cite{SGA5}; it does not identify
equal trace functions with object inverses.
SGA~6, V, \S\S2--3, with cote~133, scan~18,
supplies the actual/virtual coefficient vocabulary
\cite{SGA6,Archive}, with its existing source-access
qualification retained.
The categorical inversion lemmas and the rational-point
splitting are proved here. No noncommutative
\(\lambda\)-ring, unpublished purity theorem or
designated canonical-label identification is asserted.
